\section{Tree-growing method}
\label{sec:method}

At each internal node, conditional inference trees make two decisions. The
Stage~A feature test checks features for association with the response and
chooses a feature to split. The Stage~B threshold test then searches thresholds
only for that selected feature; if no threshold is accepted, the node becomes
terminal. In this paper, CIT denotes the resulting single-tree feature
selection method. CIF denotes a bootstrap ensemble of such trees with feature
sampling at each node and feature rankings formed by aggregating split
importances across trees.

\subsection{Stage~A and Stage~B}
\label{sec:stageA-stageB}

At a node $t$, write $(X_t,Y_t)$ for the samples reaching that node. Let
$F_{t,\mathrm{avail}}$ be the feature set carried into the node by the
tree-growing algorithm, including any feature muting inherited from ancestors.
Let $F_{t,\mathrm{nonconst}} \subseteq F_{t,\mathrm{avail}}$ be the features
that are nonconstant on the node samples. After any feature
subsampling, Stage~A considers a set
$F_t \subseteq F_{t,\mathrm{nonconst}}$; its size $m_t=|F_t|$ is the
denominator of the feature correction, the Bonferroni correction over the
Stage~A tests at the node. The benchmarked selector statistics are multiple
correlation (MC) and the randomized dependence coefficient (RDC) in
classification, and Pearson correlation (PC) and RDC in regression
(\cref{sec:experiments-methods}). Any label-independent auxiliary randomness
used by a selector statistic, such as the random projections of RDC, is denoted
by $U$ and is held fixed in the fixed-node guarantee.

When feature scanning is enabled during tree growing, Stage~A tests the
features in order of their observed selector statistic and stops at the first
rejection, so not every feature in $F_t$ is tested.
Let $E_t \subseteq F_t$ denote the evaluated subset; in the full-scan fixed-$B$
reference rule, $E_t=F_t$. The fixed-node lemma and corollaries in \cref{sec:theory} apply
to that reference rule: the features in $F_t$ and the budget $B$ are fixed
independently of the node responses, and every $p_{t,j}$ for $j\in F_t$ uses
that budget. The implemented tree-growing rule adds response-dependent node
selection, inherited feature muting, and scores from scans that stop at the
first rejection.

\paragraph{Stage~A feature test.}
In the full-scan fixed-$B$ setting, Stage~A follows the separation of variable
selection from threshold search in conditional inference trees
\citep{Hothorn2006UnbiasedRecursivePartitioning}. The original conditional
inference tests use the permutation-statistic framework of
\citet{StrasserWeber1999PermutationStatistics}; here, Monte Carlo label
permutations calibrate the configured selector statistic. Stage~A
computes a selector statistic $T^{\mathrm{sel}}_j(X_{t,j},Y_t;U)$ and
permutation $p$-value $p_{t,j}$ for every feature $j \in F_t$, then selects
%
\begin{equation}
  p_t^{\star,\mathrm{ref}} := \min_{j \in F_t} p_{t,j},
  \qquad
  M_t^\star := \{j \in F_t : p_{t,j}=p_t^{\star,\mathrm{ref}}\}.
\end{equation}
%
Let $\tau_t=\alpha_{\mathrm{sel}}/m_t$ when the feature correction is enabled and
$\tau_t=\alpha_{\mathrm{sel}}$ otherwise. If
$p_t^{\star,\mathrm{ref}} \ge \tau_t$, Stage~A returns no feature and the node
becomes terminal. If $p_t^{\star,\mathrm{ref}} < \tau_t$, the selected feature
$j_t^{\star,\mathrm{ref}}$ is drawn uniformly from $M_t^\star$, and the node
enters Stage~B.

In the implemented tree-growing rule, Stage~A uses the same decision structure
but can order features by a preliminary score and evaluate only the subset
$E_t$. The returned stored score $q_t^{\mathrm{feat}}$ is used to grow the
tree; it is not a full-scan fixed-$B$ $p$-value covered by
\cref{lem:plusone-superuniform}. Under the full-scan reference rule with the
feature correction, $\min(1,m_t p_t^{\star,\mathrm{ref}})$ is a conservative
complete-null $p$-value for the fixed-node Stage~A test. When feature muting is
enabled, features with non-rejecting Stage~A $p$-values can be removed from
descendant feature pools. When muting is disabled, descendants keep
$F_{t,\mathrm{avail}}$ except for deterministic constant feature pruning.

\paragraph{Stage~B threshold test.}
Stage~B guards against a spurious split at a node whose feature passed Stage~A:
the node splits only when some threshold of that feature separates the
responses beyond what permuted responses achieve at level $\alpha_{\mathrm{split}}$.
By default, the ctree reference implementation instead selects the split point
by maximizing a standardized two-sample statistic over the candidate splits,
without a second test of the chosen split
\citep{Hothorn2006UnbiasedRecursivePartitioning}.
Given a selected feature $j_t^\star$, Stage~B starts from midpoints between
ordered unique values of $X_{t,j_t^\star}$. The exact method uses all such
midpoints. The random method samples a bounded subset of midpoints, while the
percentile and histogram methods choose bounded representative thresholds
derived from the midpoint distribution. The histogram method with $h$ bins
returns the distinct edges of an $h$-bin histogram of the midpoints, so it
yields at most $h+1$ candidates; the 256-bin setting of the benchmark yields up
to 257. Thresholds that would leave either child below the minimum leaf size
are filtered out. If no threshold remains, Stage~B returns no split.

For the remaining finite threshold set $C_{t,j_t^\star}$, Stage~B selects the
threshold used by the tree. The tested statistic at threshold $c$ is the
weighted child impurity
\[
  \frac{|Y_t^L(c)|}{|Y_t|}\mathrm{Imp}(Y_t^L(c))
  + \frac{|Y_t^R(c)|}{|Y_t|}\mathrm{Imp}(Y_t^R(c)),
\]
where $\mathrm{Imp}$ is the chosen node impurity measure. This statistic is used
in a left-tail test, and the full-scan fixed-$B$ rule chooses the threshold with
the smallest Stage~B $p$-value, with uniform tie breaking. When threshold
scanning is enabled, thresholds are tested in order of weighted child impurity
and the search stops after the first rejecting threshold, so the returned
Stage~B score is a tree-construction score rather than a global minimum over
all retained thresholds. Let
$\tau_t^{\mathrm{split}}=\alpha_{\mathrm{split}}/|C_{t,j_t^\star}|$ when the
threshold correction, the Bonferroni correction over the candidate thresholds,
is enabled and $\tau_t^{\mathrm{split}}=\alpha_{\mathrm{split}}$ otherwise;
testing each candidate at $\tau_t^{\mathrm{split}}$ is the per-threshold rule.
If no evaluated Stage~B $p$-value is below $\tau_t^{\mathrm{split}}$, the node
becomes terminal. After a threshold passes the Stage~B test, the final minimum
required-improvement check uses the corresponding weighted impurity decrease.
Stage~B chooses the split used by the tree; the returned Stage~B score is not a
valid post-selection $p$-value. The training skeleton in \cref{alg:citree}
summarizes this flow and, for the benchmarked CIT and CIF configurations, stores
the two returned scores $q_t^{\mathrm{feat}}$ and $q_t^{\mathrm{split}}$ at each
split.

\paragraph{Permutation budgets.}
Write $\alpha_{\mathrm{adj}}$ for the level at which a single permutation test is
run: $\alpha_{\mathrm{adj}}=\alpha_{\mathrm{sel}}/m_t$ for a Stage~A test and
$\alpha_{\mathrm{adj}}=\alpha_{\mathrm{split}}/K$, with $K=|C_{t,j_t^\star}|$,
for a per-threshold Stage~B test when the corresponding correction is enabled,
and $\alpha_{\mathrm{sel}}$ or $\alpha_{\mathrm{split}}$ otherwise. The minimum
budget gives each test $B=\lceil 1/\alpha_{\mathrm{adj}}\rceil$ permutations, the
smallest budget at which a $+1$ $p$-value can fall below $\alpha_{\mathrm{adj}}$.
The auto budget, the library default, gives each test
%
\begin{equation}
  B = \max\!\left(\left\lceil \frac{1}{\alpha_{\mathrm{adj}}} \right\rceil,\
  \left\lceil \frac{z^2(1-\alpha_{\mathrm{adj}})}{\alpha_{\mathrm{adj}}} \right\rceil\right),
  \qquad z=\Phi^{-1}(1-\alpha_{\mathrm{adj}}),
  \label{eq:auto-budget}
\end{equation}
%
permutations, where $\Phi$ is the standard normal distribution function. The
benchmark uses $\alpha_{\mathrm{sel}}=\alpha_{\mathrm{split}}=0.05$ with the
minimum budget, so a Stage~A test at a node with 20 features receives 400
permutations and a per-threshold Stage~B test over 257 candidates receives
5{,}140. Adaptive stopping, the default stopping rule, checks a Beta posterior
for the exceedance probability every 32 permutations once at least
$\lceil 1/\alpha_{\mathrm{adj}}\rceil$ permutations have been drawn, and stops
when the posterior is confident on either side of $\alpha_{\mathrm{adj}}$
(\cref{prop:adaptive-size}); with no stopping, the test runs its full budget.
Under the minimum budget the first check coincides with the full budget, so
adaptive stopping never stops early.

Under the minimum budget a test rejects only when no permuted statistic is as
extreme as the observed one: $\alpha_{\mathrm{adj}}(B+1)<2$, so zero
exceedances is the only count $b$ with $(1+b)/(B+1)<\alpha_{\mathrm{adj}}$.
Every rejecting feature in Stage~A and every rejecting threshold in Stage~B
therefore returns the floor $p$-value $1/(B+1)$, and the $p$-values do not order
the passing candidates. Without scanning, the smallest-$p$ rule would draw
uniformly among them. With feature and threshold scanning, which the benchmark
uses, the first rejecting candidate in scan order is chosen. At the minimum
budget the Stage~A and Stage~B decisions are therefore an ordering by the raw
statistics, the observed selector statistic for features and the observed
weighted child impurity for thresholds, behind the Bonferroni gate formed by the
feature and threshold corrections.

\paragraph{Max-type rule for Stage~B.}
The per-threshold rule tests each candidate $c \in C_{t,j_t^\star}$ separately
at level $\alpha_{\mathrm{split}}/K$, so under the minimum budget each candidate
receives $\lceil K/\alpha_{\mathrm{split}}\rceil$ permutations and a scan of all
$K$ candidates costs on the order of $K^2/\alpha_{\mathrm{split}}$ impurity
evaluations per node. The same scaling holds in Stage~A, where each of the $m_t$
feature tests receives $\lceil m_t/\alpha_{\mathrm{sel}}\rceil$ permutations. The
implementation also provides the max-type rule, an opt-in alternative that tests
the candidates jointly. Write $S_{b,c}$ for the weighted child impurity at
candidate $c$ under the observed response ($b=0$) and under permutation
$b=1,\dots,B$. Each column is standardized by the mean $\mu_c$ and standard
deviation $\sigma_c$ of its $B+1$ entries, candidates with $\sigma_c=0$ or an
empty child are dropped, and the row statistic is
%
\begin{equation}
  T_b := \min_{c} \frac{S_{b,c}-\mu_c}{\sigma_c},
  \qquad
  p^{\max}_{t} := \frac{1+\sum_{b=1}^{B}\ind\{T_b \le T_0\}}{B+1},
  \label{eq:maxt-stat}
\end{equation}
%
with the split placed at $\arg\min_c (S_{0,c}-\mu_c)/\sigma_c$, the standardized
rather than the raw impurity minimum, so the two rules can accept different
thresholds on the same feature. Stage~B rejects when
$p^{\max}_t<\alpha_{\mathrm{split}}$, without division by $K$. When every
candidate is dropped, the test returns $p^{\max}_t=1$ and no split.
\Cref{cor:stageB-maxt} gives the fixed-node guarantee.

The max-type rule is a single-step max-T statistic in the sense of
\citet{WestfallYoung1993ResamplingBasedMultipleTesting}, applied to a minimum
because smaller impurity is more extreme, with the empirical null
standardization of \citet{PollardVanDerLaan2004NullDistribution}. It has the
same standardize-then-optimize form as the conditional inference split
statistic of \citet{Hothorn2006UnbiasedRecursivePartitioning}, with Monte Carlo
column moments replacing the exact conditional moments of
\citet{StrasserWeber1999PermutationStatistics} because impurity is not a linear
statistic. It is a maximally selected statistic in the sense of
\citet{LausenSchumacher1992MaximallySelectedRankStatistics} and
\citet{HothornZeileis2008GeneralizedMaximallySelectedStatistics}, which is what
ctree's split search computes exactly for linear statistics. Our additions are
the Monte Carlo standardization, which extends that construction to the
impurity criterion, and the comparison of its cost with that of the
per-threshold rule. Standardization is intended to keep candidates near the
tails of $X_{t,j_t^\star}$, which isolate few observations and have wide null
distributions, from dominating the minimum on heavy-tailed responses; we did not
run an unstandardized arm to test this.

Validity of the max-type rule needs only about $1/\alpha_{\mathrm{split}}$
permutations, and the implementation uses $B=999$ under every budget rule, so
testing all $K$ candidates costs on the order of $BK$ impurity evaluations
against $K^2/\alpha_{\mathrm{split}}$ for the per-threshold rule under the
minimum budget. These counts are the worst case in which every candidate is
tested. With threshold scanning, the realized cost of the per-threshold rule is
the number of candidates scanned times $\lceil K/\alpha_{\mathrm{split}}\rceil$.
The max-type rule needs fewer permutations per node only when
$K > \alpha_{\mathrm{split}} B \approx 50$: at 16 bins, which yield 17
candidates, the per-threshold rule needs $\lceil 17/0.05\rceil = 340$
permutations per candidate against the 999 of the max-type rule. A
per-threshold rule held at a fixed $B$ also costs $BK$, but it cannot reject
once $\alpha_{\mathrm{split}}/K < 1/(B+1)$, which at $B=999$ is every node with
more than about 50 candidates. The benchmark configurations use the
per-threshold rule, the library default, and \cref{sec:results-stageB} reports
the measured calibration and cost of both rules and the max-type rerun of the
CIT and CIF configurations.

\subsection{Monte Carlo permutation \texorpdfstring{$p$}{p}-values}
\label{sec:plusone}

Let $T_{\mathrm{obs}}$ be the observed statistic and $T_1,\dots,T_B$ the same statistic
computed on $B$ random label permutations. For a right tail test we use the
$+1$ $p$-value of Phipson and Smyth
\citep{NorthCurtisSham2002EmpiricalPValuesMonteCarlo,PhipsonSmyth2010PermutationPValues},
where $\ind\{\cdot\}$ denotes the indicator function:
%
\begin{equation}
  p_{\mathrm{MC}}^{\ge} := \frac{1 + \sum_{b=1}^B \ind\{T_b \ge T_{\mathrm{obs}}\}}{B+1}.
\end{equation}
%
For a left tail test we analogously define
%
\begin{equation}
  p_{\mathrm{MC}}^{\le} := \frac{1 + \sum_{b=1}^B \ind\{T_b \le T_{\mathrm{obs}}\}}{B+1}.
\end{equation}
%
Stage~A uses the right tail convention and Stage~B the left tail convention;
ties count as exceedances in both cases.

\begin{algorithm}[H]
  \caption{Conditional inference tree training skeleton (Stage~A $\rightarrow$ Stage~B)}
  \label{alg:citree}
  \begin{algorithmic}[1]
    \Require Training data $(X,Y)$, hyperparameters $\theta$
    \Ensure Trained tree $\mathcal{T}$
    \State $p \gets$ number of columns of $X$
    \State \Return \Call{GrowNode}{$X,Y,1,\{1,\dots,p\},\theta$}
    \Function{GrowNode}{$X_t,Y_t,d,F_{\mathrm{avail}},\theta$}
      \If{\Call{Stop}{$X_t,Y_t,d,\theta$}}
        \State \Return \Call{Leaf}{$Y_t$}
      \EndIf
      \State $(j_t^\star, q_t^{\mathrm{feat}}, m_t, F'_{\mathrm{avail}}) \gets \Call{StageA}{X_t,Y_t,F_{\mathrm{avail}},\theta}$
      \If{$j_t^\star = \bot$}
        \State \Return \Call{Leaf}{$Y_t$}
      \EndIf
      \State $(c_t^\star, q_t^{\mathrm{split}}) \gets \Call{StageB}{X_t,Y_t,j_t^\star,\theta}$
      \If{$c_t^\star = \bot$}
        \State \Return \Call{Leaf}{$Y_t$}
      \EndIf
      \State $(X_t^L,Y_t^L,X_t^R,Y_t^R) \gets$ \Call{Partition}{$X_t,Y_t,j_t^\star,c_t^\star$}
      \State $\Delta_t \gets \Call{ImpurityDecrease}{Y_t,Y_t^L,Y_t^R,\theta}$
      \If{$\Delta_t$ is below the minimum required impurity decrease}
        \State \Return \Call{Leaf}{$Y_t$}
      \EndIf
      \State Add $\Delta_t$ to the feature-importance accumulator for $j_t^\star$
      \State $u \gets$ internal node with split $(j_t^\star,c_t^\star)$ and stored scores $(q_t^{\mathrm{feat}},q_t^{\mathrm{split}})$
      \State $u.\mathrm{left} \gets \Call{GrowNode}{X_t^L,Y_t^L,d+1,F'_{\mathrm{avail}},\theta}$
      \State $u.\mathrm{right} \gets \Call{GrowNode}{X_t^R,Y_t^R,d+1,F'_{\mathrm{avail}},\theta}$
      \State \Return $u$
    \EndFunction
  \end{algorithmic}
\end{algorithm}

\subsection{Forests and runtime hyperparameters}
\label{sec:bootstrap}

CIF combines bootstrap trees and averages their predictions. Each CIF tree uses
the same Stage~A and Stage~B node rules as CIT
\citep{Breiman2001RandomForests,HothornBuehlmannDudoitMolinaroVanDerLaan2006SurvivalEnsembles}.
For feature ranking, CIF sums each fitted tree's normalized split importances
and renormalizes the forest vector (\cref{sec:rank-construction}). CIF-all
disables feature subsampling: every nonconstant available feature can enter
Stage~A at each node. This changes the Stage~A feature set, the feature
correction, and the per-node computational cost. The fixed-node results of
\cref{sec:theory} apply to the node tests of each tree with the scope stated in
\cref{app:theory-scope}; that scope covers subsampled and unsampled fits, and not
bootstrap samples, whose duplicated rows break the exchangeability the results
assume.

Across the CIT and CIF benchmark configurations, we hold fixed the feature and
threshold corrections, the per-threshold rule, the minimum budget with adaptive
stopping (which therefore never stops early), feature muting, feature scanning,
threshold scanning, the histogram method with 256 bins, and bootstrap sampling
for CIF. The settings that vary are the selector statistic and honesty. Honesty
grows the structure of each tree on one part of its training sample and
re-estimates the leaf values on the held-out remainder
\citep{WagerAthey2018CausalForests}; the benchmark uses an even split when
honesty is on.

Runtime depends on the realized work at each node: features tested in Stage~A,
thresholds generated and retained in Stage~B, permutations evaluated before
stopping, and features carried into the node. Adaptive stopping can reduce the
number of permutations evaluated for each test when the budget exceeds the
minimum budget
\citep{BesagClifford1991SequentialMonteCarloPValues,Gandy2009SequentialImplementationMonteCarloTests}.
The histogram method reduces the Stage~B candidate set from all midpoints to at
most $h+1$ bin edges. Feature scanning orders features by the observed selector
statistic and threshold scanning orders thresholds by the observed weighted
child impurity; both stop at the first rejection, independently of the stopping
rule inside each permutation test. Feature muting can reduce the feature pool
inherited by descendants.
Appendix~\ref{app:methods} gives algorithm and
reproducibility details, Appendix~\ref{app:complexity-details} gives runtime and
memory accounting, and \cref{sec:results-practical} reports the measured timing
effects.

\subsection{Feature rankings}
\label{sec:rank-construction}

The experiments evaluate each method through the ranking it produces on the
training fold. CIT and CIF rankings are fitted split-importance rankings, not
rankings of Stage~A $p$-values. For a fitted CIT, the importance of feature $j$
is the sum, over the internal nodes $t$ whose accepted split is on $j$, of the
impurity decrease
\[
  \mathrm{Imp}(Y_t) - \frac{|Y_t^L|}{|Y_t|}\mathrm{Imp}(Y_t^L)
  - \frac{|Y_t^R|}{|Y_t|}\mathrm{Imp}(Y_t^R),
\]
the parent impurity minus the weighted child impurity at the accepted threshold,
with no further weight for the node's share $|Y_t|/n$ of the sample. If the
total importance is positive, the vector is normalized to sum to one; if the
tree contributes no positive impurity decrease, the vector remains zero. CIF
sums these tree importance vectors over trees and then normalizes the forest
vector in the same way. Stage~A determines which features enter threshold
optimization, and the importance scores only the accepted splits by their local
reduction in impurity.

This unweighted sum differs from scikit-learn's impurity importance, which
multiplies each decrease by the node share so that root splits dominate. It
counts a deep split on a small node as much as a root split, and because every
accepted split has passed both the Stage~A and the Stage~B test, the ranking
behaves like a gain-weighted count of accepted splits; the ranking ablation
(\cref{tab:cif-ranking-ablation}) compares it with a split-count ranking. Larger
importances are ranked earlier. Features never used by a split receive zero
importance, so beyond a tree's support the order of the ranking is the tie
order of the implementation, which leaves the zero-importance features in
roughly ascending column order. Every synthetic generator shuffles its planted
columns, so that order carries no information about which features are planted;
the high-cardinality and correlated-noise designs append their noise columns
after the others, so those columns fall toward the end of the tail.

This split-importance ranking differs from standard random-forest impurity
importance in how a feature enters threshold search. In CART-style trees, the
feature and threshold are chosen together by an impurity search, so a feature
with more possible thresholds has more chances to produce a large apparent
impurity decrease \citep{Strobl2007BiasRFVariableImportance}. In CIT and CIF,
Stage~A first selects a feature by a permutation test; Stage~B then searches
thresholds only for that feature. This is designed to reduce the advantage that
high-cardinality features get from having more possible thresholds. The
threshold correction of the per-threshold rule works in the opposite direction:
more candidates mean a stricter level $\alpha_{\mathrm{split}}/K$ for each
candidate and a lower probability of any split, a reverse cardinality bias
against features with many candidate thresholds. The max-type rule tests the
candidates jointly at $\alpha_{\mathrm{split}}$, so its level does not depend on
$K$; \cref{sec:results-stageB} and \cref{tab:weak-signal-cardinality} report
how each rule behaves. We use these fitted split-importance scores as rankings
and evaluate them through downstream top-$k$ prediction and synthetic feature
recovery.
